% Generated draft: requires mathematical and source review.
\documentclass[11pt]{article}
\usepackage[margin=25mm]{geometry}
\usepackage{fontspec}
\setmainfont{Latin Modern Roman}
\usepackage{amsmath,amssymb,mathtools,unicode-math}
\setmathfont{Latin Modern Math}
\usepackage{enumitem,needspace,xurl,hyperref}
\setlength{\parindent}{0pt}
\setlength{\parskip}{5pt}
\setlength{\emergencystretch}{3em}
\newcommand{\N}{\mathbb{N}}
\newcommand{\Z}{\mathbb{Z}}
\newcommand{\Q}{\mathbb{Q}}
\newcommand{\R}{\mathbb{R}}
\newcommand{\C}{\mathbb{C}}
\newcommand{\F}{\mathbb{F}}
\newcommand{\A}{\mathbb{A}}
\newcommand{\PP}{\mathbb{P}}
\newcommand{\E}{\mathbb{E}}
\newcommand{\eps}{\varepsilon}
\newcommand{\dd}{\,\mathrm{d}}
\begin{document}
\section*{UnsolvedMath ID: 10600015}
\textbf{Source problem code:} catalogue entry
\subsection{Definitions and mathematical statement}
A virtual knot can be represented by a knot diagram with classical over/under crossings and virtual crossings, modulo the classical Reidemeister and virtual detour moves; equivalently one uses knots in thickened surfaces with the appropriate stabilization equivalence. A common alternating-diagram convention requires successive encounters with classical crossings to alternate over and under, ignoring virtual crossings. Whether this is the precise definition sought by the source must be checked.

The compound problem asks:
\begin{enumerate}
\item Define and classify alternating virtual knots. A diagram-level alternating condition must be distinguished from the property of a virtual knot admitting such a representative, and the requested classification must specify its equivalence.
\item Formulate and prove a virtual analogue of the Tait flyping conjecture. A classical flype rotates an adjacent tangle through a crossing. The intended virtual/surface moves, reducedness or primeness hypotheses, and handling of supporting surfaces must be specified; the excerpt does not justify an unrestricted claim for every pair of alternating virtual diagrams.
\item Classify all alternating weaves on surfaces \emph{without stabilization}. Thus adding or removing surface handles is not an allowed identification in this part. The source does not define weave, impose cellularity, compactness or periodicity, or say whether the fixed surface is marked or considered up to homeomorphism.
\end{enumerate}
Surface crossing-decorated graphs are a natural setting for the final part, but their exact admissibility and equivalence are missing. It must not be collapsed into stable virtual-knot classification.

The summary reports a flyping analogue solved by Kindred (2022), related alternating-link extensions, and work on surface weaves in 2020, while retaining full classification as a programme. Exact references and hypotheses are absent, so these are reported progress claims rather than verified resolutions of all three tasks. The source's ``Compare'' citation was lost in extraction.
\subsection{Short English statement}
Define and classify alternating virtual knots, formulate/prove the proper flyping analogue, and separately classify alternating weaves on fixed surfaces without stabilization. The reported flyping progress and missing comparison citation require verification.
\subsection{Sources}
\begin{itemize}
\item Local UnsolvedMath registry, permanent ID 10600015, source code catalogue entry; no proposer or proposal year supplied; updated 2026-08-06.
\item \emph{Unsolved Problems in Virtual Knot Theory and Combinatorial Knot Theory}, revised survey problem-list item 15: \url{https://arxiv.org/abs/1409.2823}. Registry access date 2026-07-29; source-TeX extraction. Provenance reports 60 top-level items, with compound items retained. Present primary-literature status and redistribution terms require review.
\item The comparison citation after ``Compare'' is absent. The supplied summary mentions Kindred (2022) and unnamed 2020 weave work without full bibliographic entries or theorem hypotheses; these were not independently checked.
\end{itemize}
\end{document}
